\documentclass{article}
\usepackage[T1]{fontenc}
\usepackage{microtype,graphicx,booktabs,amsmath,amssymb,mathtools,array,tabularx,multirow}
\usepackage{xcolor,url}
\newcommand{\theHalgorithm}{\arabic{algorithm}}
\usepackage{mlsys2025}
\usepackage{hyperref}
\usepackage[all]{hypcap}
\usepackage{tikz,pgfplots}
\usetikzlibrary{arrows.meta,positioning,fit,backgrounds,calc}
\usepgfplotslibrary{groupplots}
\pgfplotsset{compat=1.18}
\hypersetup{colorlinks=true,linkcolor=black,citecolor=black,urlcolor=black,
  pdfauthor={Anonymous Authors},pdftitle={EviCastKV: Shared-Source Storage for Multimodal KV Caches}}
\mlsystitlerunning{EviCastKV: Shared-Source Storage for Multimodal KV Caches}
\mlsyskeywords{multimodal inference, KV cache, compression, GPU memory, reproducibility}
\definecolor{eviblue}{HTML}{1565A8}
\definecolor{exactgreen}{HTML}{227751}
\definecolor{patchorange}{HTML}{A85D10}
\definecolor{riskred}{HTML}{B13C3C}
\definecolor{muted}{HTML}{59636C}
\newcommand{\Evi}{EviCastKV-S}
\newcommand{\Full}{FullKV}
\newcommand{\Top}{Prompt Top-K}
\newcommand{\MiB}{\,\mathrm{MiB}}
\newcommand{\softmax}{\operatorname{softmax}}
\newcommand{\bytes}{\operatorname{bytes}}
\newcommand{\Rope}{\operatorname{RoPE}}
\newcommand{\Approx}{\operatorname{Approx}}
\newcommand{\Restore}{\operatorname{Restore}}
\newcommand{\SourceRev}{\texttt{a6f689f}}
\newcommand{\smallpath}[1]{\nolinkurl{#1}}
\newcolumntype{Y}{>{\raggedright\arraybackslash}X}
\newcommand{\EviInitial}{41.14}
\newcommand{\EviPeak}{83.52}
\newcommand{\EviMaxPeak}{84.24}
\newcommand{\EviFinal}{73.22}
\newcommand{\EviCider}{119.77}
\newcommand{\EviMedian}{6.55}
\newcommand{\EviTail}{6.72}
\newcommand{\FullCider}{116.83}
\newcommand{\FullInitialReduction}{86.35}
\newcommand{\FullFinalReduction}{79.94}
\newcommand{\FullPeakReduction}{77.47}
\newcommand{\TopInitialReduction}{18.89}
\newcommand{\TopFinalReduction}{11.57}
\newcommand{\TopPeakReduction}{10.03}
\newcommand{\LookFaster}{86.14}
\newcommand{\PrefixFaster}{11.07}
\newcommand{\TopFaster}{12.35}
\newcommand{\VLFaster}{22.13}


\begin{document}
\twocolumn[
\mlsystitle{\texorpdfstring{EviCastKV: Shared-Source Storage for Multimodal KV Caches\\Supplementary Material}{EviCastKV: Supplementary Material}}
\begin{mlsysauthorlist}\mlsysauthor{Anonymous Authors}{anon}\end{mlsysauthorlist}
\mlsysaffiliation{anon}{Anonymous Institution}
\mlsyscorrespondingauthor{Anonymous Authors}{anonymous@example.invalid}
\vskip .2in
]
\printAffiliationsAndNotice{}
\appendix
\section{Reading guide and claim boundary}
This supplement expands the representation, execution and measurement
contracts of the accompanying paper. The evaluated system is \Evi{}, the
shared-source precision/storage variant. The suffix S distinguishes its
scope; it does not introduce a second measured competitor. Its legacy
machine-readable method identifier is \texttt{mosaickv\_exact}. That identifier
must not be read as evidence that all prompt K/V are bitwise exact or that
the original six-component EviCastKV/MosaicKV compressor ran.

The experiment establishes completed operating points on an inspected
256-image COCO cohort. It does not establish an independent confirmation,
a physical-byte-matched selection-policy comparison, an equal-quality
memory frontier, a whole-request peak certificate, or performance of
inactive mechanisms. A coherent interpretation of the main paper requires
all of these restrictions; they are not exceptions hidden by an average.

\section{Executed configuration}
\begin{table*}[t]
\centering\small% Exported from hash-verified measured records; do not hand edit.
\begin{tabular}{lp{0.64\linewidth}}
\toprule
Field & Executed setting\\
\midrule
Model revision & \nolinkurl{b234b804b114d9e37bb655e11cbbb5f5e971b7a9}\\
COCO revision & \nolinkurl{c3213b54e173ae86517d481de17f31439d6f8c72}\\
Producer revision & \nolinkurl{a6f689fcdfae2b9023045aa412c6d90d8ce79012}\\
Python & 3.11.15\\
PyTorch / CUDA runtime & 2.8.0 / 12.8\\
Transformers / Accelerate & 4.55.2 / 1.13.0\\
lmms-eval / pycocoevalcap & 0.7.2 / 1.2\\
GPU / driver & NVIDIA A100-SXM4-80GB / 595.71.05\\
Requests / output policy & Batch one / 128 predictions, 127 append forwards\\
Processor / dataset split & Same model revision / val\\
Seed / attention & 20261106 / eager\\
Installed, not exercised & vLLM 0.10.2; SGLang 0.5.1.post3\\
\bottomrule
\end{tabular}

\caption{Executed configuration and complete immutable revisions. The source
revision belongs to the measured producer, not the later paper-writing
worktree. Installed serving packages are explicitly not exercised.}
\label{tab:environment}
\end{table*}
Table~\ref{tab:environment} records the executed environment. The workload
uses one image and one request per GPU. The checkpoint has 32 layers, 32
KV heads and head dimension 128. All weights are BF16. Original decoder
K/V are BF16; the prompt source has its separately declared three-bit
representation. Mixed precision at the representation boundary must not be
described as a lower-precision model.

Every worker uses one A100-SXM4-80GB and 12 CPU cores. The campaign admits
at most six workers concurrently, not six GPUs per request. Its execution
passes ran separately. GPU sharing on the node is recorded; clocks and
power state are observations, not proof of exclusive node ownership or
fixed-frequency execution. The deterministic setup sets
\texttt{CUBLAS\_WORKSPACE\_CONFIG=:4096:8} before Python, records the seed and
uses greedy token selection. Determinism checks remain empirical under
the pinned software rather than a promise across arbitrary CUDA versions.

The exact prompt is: ``Provide a one-sentence caption for the provided
image.'' The common processor and chat template produce an unpadded
603-position decoder prefix. Every method computes 128 predictions. The
first comes from original prefill logits; predictions 2--128 require 127
additional decoder forwards. Consequently 128 predictions are not 128
new cache positions.

\section{Representation and storage layout}
\subsection{Why a common source can reduce persistence}
In a multi-head-attention layer, both key and value projections consume
the same normalized hidden input $X_\ell\in\mathbb R^{n\times d}$.
Their dense BF16 outputs occupy $2nd$ words, whereas a single source
requires $nd$ words before compression. Reusing this algebraic input is
not a new general principle: latent-cache and low-rank projection systems
also exploit reconstruction opportunities \citep{chang2024palu}. The
specific questions here are whether the measured format can retain all
prompt positions, protect selected exact entries, encode the generated
history, and reconstruct fast enough under an explicit ownership boundary.

For grouped-query attention, $d$ can be much larger than the combined K/V
width. The same dimensional argument can then become unfavorable. The
present measurements do not establish a useful GQA format, Qwen support,
or cross-family portability \citep{ainslie2023gqa}.

\subsection{Prompt-source format}
The prompt packet contains grouped affine three-bit source codes, scale
and offset metadata, exact protected K/V payloads, positional information,
numerical patches, index metadata and applicable lossless code transforms.
Groups have 128 elements along the configured source axis. The constant
group case uses a guarded scale rather than division by zero. Padding and
allocation slack remain part of the backing storage, even when no logical
element reads them.

The source approximation affects unprotected entries across the full
prompt. Exact overlays overwrite only their protected head-position
locations. Globally exact rows can be excluded from the lossy payload to
avoid duplicate persistence; this does not remove their logical positions.
The representation is therefore neither a pure token subset nor a
prototype-merged cache. No averaging across incompatible rotary phases is
performed: original phases are retained and applied during reconstruction
\citep{su2021rope}.

\subsection{Protection rule and its limits}
The frozen protection rule combines a head-local mean over the last 16
prompt-query positions with a shared prompt prior, at weights 0.5 and
0.5. The prior is formed from original prompt attention, not generated
answers or benchmark references. Deterministic ranking and the registered
source-slot allocation produce a selected set; the last prompt position
is mandatory. The first four and last eight selected positions and
selected text positions form head-local exact regions.

This rule does not guarantee that every answer-critical visual or OCR
entry is protected. Nor is the historical half-source-slot convention
equivalent to half the physical bytes. Different exact overlays,
adaptive layer quotas, merge payloads and metadata explain why the six
methods in this study do not produce equal encoded caches.

\subsection{Two distinct numerical references}
Let $X$ be the original source and $\widehat X$ its decoded approximation.
There are two different comparisons:
\begin{enumerate}
\item \emph{Quantization/task distortion}: compare computation from
      $\widehat X$ and protected overlays with unrestricted computation
      from $X$. This is a lossy modeling change; caption quality must be
      measured. Numerical patches do not eliminate this distortion.
\item \emph{Transport distortion}: compare a custom reconstruction
      kernel with native frozen-weight projection of the same
      $\widehat X$, the same overlays and the same original rotary phases.
      This checks the arithmetic implementation of a chosen representation.
\end{enumerate}
The registered prompt transport envelope uses relative tolerance 0.01 and
absolute tolerance $2^{-6}$ under the pinned dtype predicate. Sparse
native-reference words repair violations that the format explicitly
permits. Protected-row mismatches are rejected. This is construction-time
numerical handling, not uncertainty-guided decode repair, learned value
fitting or benchmark-label calibration.

The generated-history path has a stricter, bitwise BF16 K/V contract.
Source and correction coding can be lossless even though the prompt that
caused those generated forward passes was approximate. These two
statements are consistent and must not be collapsed into a claim that
the entire request equals FullKV.

\subsection{History sealing and unfinished tails}
Each generated token contributes its actual normalized input and original
logical position. The current path seals 16-token chunks, retains the
unfinished tail, and preserves continuity when consolidating archived
chunks. Native projections and original rotary phases reconstruct the
history. Sparse exact-word corrections are checked through BF16 word
identity. Lossless payload coding falls back to raw storage when coding
would expand it.

There is no generated-token history eviction in this study. Capacity is
not kept constant by discarding old answers. An append-only logical cache
can nevertheless have compressed physical backing. The final and peak
measurements include the accumulated history; reporting the initial
prompt packet alone would overstate the lifecycle benefit.

\section{Request lifecycle and cache ownership}
\subsection{Construction and publication}
The following lifecycle describes the implementation, not a second model
prefill or an offline preprocessing stage:
\begin{enumerate}
\item Enter the ordinary multimodal request and process its image/text.
\item During the one original prefill, capture the layer source, original
      K/V, prompt attention needed by that policy, and rotary metadata.
\item Construct a private encoded packet; validate shape, finite values,
      protected entries and the numerical predicate before publication.
\item Retire superseded dense cache owners. Bind the packet and charge
      the live metadata introduced by binding.
\item At each generated forward, reconstruct the current layer into a
      bounded lease, consume it through attention and retire or reuse
      only request-private transient storage.
\item Append the actual generated-history source at its original logical
      next position; seal chunks under the declared lossless contract.
\item Export logits, derive the EOS-aware caption, retire private replay
      state and return. End-to-end timing includes these steps.
\end{enumerate}

Publication is a state transaction. A serializable packet size is not the
same as initial bound-cache storage, and a tensor view's logical size is
not the size of its backing allocation. No uncharged full-cache owner may
remain resident merely to make a later correction easier.

\subsection{First-token boundary}
The evaluated boundary uses the original complete prompt prefill. Its
last logits determine the first generated token, then the encoded cache
supports subsequent forwards. All six methods therefore share their
first predicted token for each image; the exporter checks this invariant.
That token is not evidence of approximate-cache output parity.

A different historical diagnostic withheld the final prompt token and
then evaluated it against a compressed prefix. That boundary needs its
own FullKV numerical control, observation mask and no-double-feed checks.
It is not silently substituted here, and earlier failures at that boundary
are not resolved by the present boundary's successful execution. No
streamed first-token emission was measured; prefill-logit availability
must not be labeled a measured user-visible TTFT.

\subsection{Conditional uniform-layout replay}
The acceleration predicate requires the supported 32-layer configuration,
uniform layer prompt lengths and the bounded prompt-plus-update extent
of the current eager core. For the present 603-position prompt and 127
updates, that extent is below its 1,024-position limit. Eligible requests
use private split execution graphs with explicit dtype handling. Ineligible
adaptive layouts use their ordinary eager path without quota changes or
padding solely to qualify for replay.

The common capability was requested for all compressed methods, but all
256 \Evi{} layouts qualified and all four frozen baseline layouts failed
the shape predicate. This is a material implementation difference. The
latency result is for complete representation/layout/runtime operating
points; it is not a causal estimate of the source codec alone. Graph setup,
capture, retirement and charged workspace remain inside the request.
CUDA execution graphs do not constitute the historical semantic evidence
graph, and their activity cannot validate that inactive algorithm component.

\section{Measurement definitions and accounting exclusions}
\subsection{Persistent and decode-owned storage}
For live tensors $\mathcal T(t)$ at a declared boundary, define
\begin{equation}
 B(t)=\sum_{s\in\mathrm{unique\_backings}(\mathcal T(t))}
       \mathrm{allocation\_bytes}(s).
\end{equation}
Deduplication uses storage identity, not tensor-object identity. A view
does not create a second charged owner, and a small view does not shrink
the charge of a larger retained backing. Codes, headers, exact overlays,
patches, stored positions and relevant metadata are included.

Initial storage is measured after construction and binding; final storage
after the prescribed generated forwards. The audited decode-owned peak
observes reconstruction leases and other live cache tensors, including
native append overlap where relevant. Means of per-request peaks and the
maximum observed request peak are different statistics and both are
reported.

The audit does not certify a whole-request KV peak spanning every prefill
and construction transient, nor total process GPU memory. It excludes
model weights, general attention activations, vendor-internal resources
and unrelated process allocations. CUDA allocator reserved memory is not
the same quantity as unique live cache backing. Memory percentages in
the paper must not be translated into an identical percentage of GPU HBM
capacity or a proven serving-batch increase.

\begin{table*}[t]
\centering\small% Exported from hash-verified measured records; do not hand edit.
\begin{tabular}{lrrr}
\toprule
Reference point & Initial & Mean decode peak & Final\\
\midrule
FullKV reference & 86.35\% & 77.47\% & 79.94\%\\
Prompt Top-K (A) & 18.89\% & 10.03\% & 11.57\%\\
PrefixKV (R+A) & 9.97\% & 2.90\% & 5.85\%\\
VL-Cache (R+A) & 22.82\% & 11.52\% & 14.25\%\\
LOOK-M (R+A) & 18.88\% & 10.03\% & 11.57\%\\
\bottomrule
\end{tabular}

\caption{Reduction of \Evi{} relative to each measured operating point.
All percentages are computed from measured cohort means. These are
different-byte comparisons, not matched-budget or equal-quality results.}
\label{tab:reductions}
\end{table*}

\subsection{Full-request and phase time}
The headline latency is synchronized completed-request wall time, measured
without the storage observer. It includes actual media handling,
preprocessing, prefill, construction, graph setup, decode, CPU logits
export, EOS-aware detokenization and private-state retirement. Model
loading, writing artifacts and evaluating CIDEr are excluded. A request
does not reuse another timed request's prefill or captured graphs.

\begin{table*}[t]
\centering\small\setlength{\tabcolsep}{4pt}
% Exported from hash-verified measured records; do not hand edit.
\begin{tabular}{lrrrrr}
\toprule
Method & Prepare & Prefill/construct/retire & Decode & Full request & Replay eligible\\
\midrule
\textbf{EviCastKV-S} & 0.011 & 0.802 & 5.701 & 6.568 & 256/256\\
Prompt Top-K (A) & 0.011 & 0.799 & 6.642 & 7.508 & 0/256\\
PrefixKV (R+A) & 0.011 & 0.688 & 6.641 & 7.395 & 0/256\\
VL-Cache (R+A) & 0.011 & 5.336 & 7.911 & 13.349 & 0/256\\
LOOK-M (R+A) & 0.011 & 42.051 & 6.640 & 48.742 & 0/256\\
\bottomrule
\end{tabular}

\caption{Mean recorded phase wall times (seconds), over the uninstrumented
256-image pass. Prefill/construct/retire is a combined region, not a
compression-only measurement. The full request contains additional
operations and is not defined by summing these columns. Replay eligibility
is actual engagement under the shared predicate.}
\label{tab:phases}
\end{table*}
The phase table is deliberately limited to the compressed paths, whose
recorded keys define compatible regions. FullKV uses a different diagnostic
phase dictionary; absent regions are not filled with invented zeros.
Inclusive host and device intervals can overlap. In particular, a CUDA
stream interval containing gaps is not exclusively GPU kernel time, and
its sum with nested host durations would double-count work
\citep{pytorch2025cuda}.

\subsection{Correctness checks are not benchmark wins}
The canonical CPU suite records 4,548 passing tests and two deselected GPU
tests. GPU execution then validates six audit/non-audit canaries, with all
128 exported logits bitwise equal for each method's instrumented and fast
path. The complete memory and timing passes additionally have identical
token IDs and storage endpoints for every method--image pair. The repeated
study preserves those tokens and endpoints as well.

These checks rule out the measured instrumentation changing outputs, and
help detect state or accounting drift. They do not prove that lossy
prompt-source K/V equal FullKV, that another backend is supported, that
all possible requests are isolated, or that every numerical threshold in
the original six-component design passed. At retention one, the declared
identity path must bypass lossy conversion; that general contract is not
replaced by the audit/non-audit comparison.

\section{Captions, references and denominators}
\begin{table*}[t]
\centering\small% Exported from hash-verified measured records; do not hand edit.
\begin{tabular}{lrrr}
\toprule
Method & Empty & Limit reached & Mean scored tokens\\
\midrule
FullKV reference & 0 & 0 & 14.27\\
\textbf{EviCastKV-S} & 0 & 0 & 14.18\\
Prompt Top-K (A) & 0 & 0 & 14.43\\
PrefixKV (R+A) & 0 & 42 & 30.98\\
VL-Cache (R+A) & 0 & 0 & 14.07\\
LOOK-M (R+A) & 0 & 0 & 14.00\\
\bottomrule
\end{tabular}

\caption{Caption parsing over all 256 images per method. A limit-hit caption
remains in the metric denominator and is not an execution failure silently
discarded. The fixed-work decode still computes 128 predictions after an
early EOS; scoring uses the EOS-aware prefix.}
\label{tab:captions}
\end{table*}
Every row retains its generated token IDs, caption, reference list, terminal
metadata, timing and storage endpoints. Some COCO records provide more
than five captions; all supplied references are accepted. Empty captions
would be scored and recorded, not removed. The accepted study has none.

The PrefixKV precision/storage adaptation reaches the prediction limit in
42 of 256 captions; all five other methods do so in zero cases.
Its observed caption length and quality therefore cannot establish that
official PrefixKV fails at its published VQA or PPL operating point. No
PrefixKV teacher-forced PPL or native FullKV-reference ROUGE-L is measured
by this COCO experiment. Likewise, CIDEr cannot substitute for MileBench
Accuracy/ROUGE-L or VL-Cache's other paper-native routes.

\section{Statistical procedures}
\subsection{Corpus CIDEr is not a mean of shard scores}
Let $I$ denote an image multiset, $y^m_i$ method $m$'s caption and
$R_i$ its references. Write the complete evaluator as
\begin{equation}
 C_m(I)=\mathcal C\bigl(\{(y^m_i,R_i):i\in I\}\bigr).
\end{equation}
The reducer includes reference document frequencies and TF--IDF weights
\citep{vedantam2015cider}. A bootstrap draw $I_b$ resamples 256 paired
image indices with replacement, including multiplicities. It recomputes
the corpus quantities on $I_b$ and forms
$D_b=C_{\mathrm{Evi}}(I_b)-C_m(I_b)$. Averaging precomputed shard scores
or original-corpus per-image CIDEr values is not this procedure.

The point estimate uses the original complete cohort. We use 10,000 draws
and seed 20261106. The unadjusted interval uses the 0.025 and 0.975
quantiles. For the five registered quality contrasts, Bonferroni intervals
use tail probability $0.05/(2\cdot5)$ per contrast, producing a familywise
95\% interval family \citep{efron1979bootstrap}.

Positive point estimates do not establish superiority when the lower
bound remains negative. Intervals crossing zero do not establish
equivalence. A zero-margin preservation claim needs its own justified
test and a nonnegative lower bound under the stated multiplicity control.
The present inspected cohort cannot be relabeled untouched confirmation
by a different statistical reducer.

\subsection{Repeated request-time ratios}
The separate repeated experiment contains six preregistered image indices,
20 repetitions per compressed method and five warmups. It has 600 measured
rows. Methods are ordered under the registered randomized schedule, and
model loading remains outside the timer. For each reference adaptation,
the ratio is mean \Evi{} request time divided by mean baseline request
time. Nested paired resampling retains the image/repetition structure;
four contrast intervals are familywise adjusted.

\begin{table*}[t]
\centering\small% Exported from hash-verified measured records; do not hand edit.
\begin{tabular}{lrrr}
\toprule
Baseline & Time ratio & Adjusted 95\% interval & Lower time\\
\midrule
LOOK-M (R+A) & 0.1386 & [0.1325, 0.1451] & 86.14\%\\
PrefixKV (R+A) & 0.8893 & [0.8840, 0.8936] & 11.07\%\\
Prompt Top-K (A) & 0.8765 & [0.8717, 0.8807] & 12.35\%\\
VL-Cache (R+A) & 0.7787 & [0.7697, 0.7850] & 22.13\%\\
\bottomrule
\end{tabular}

\caption{Separate six-image repeated study. ``Lower time'' is $100(1-r)$,
where $r$ is the ratio of means, not the mean of per-request ratios. FullKV
is excluded from compressed-method latency inference.}
\label{tab:repeated}
\end{table*}
The 256-image p50/p95 describe a broader distribution with one timed
observation per method--image pair. They are not 20-repeat quantiles for
all 256 images, and are not pooled with the six-image repeated study.
This distinction particularly matters when local construction cost varies
substantially across images.

\section{Active and inactive design components}
\begin{table*}[t]
\centering\small
\begin{tabularx}{\textwidth}{l c Y}
\toprule
Historical mechanism & Active? & Interpretation of this study\\
\midrule
Future-query forecasting & No & Prompt protection is not a temporary future rollout.\\
Sparse cross-modal evidence graph & No & CUDA replay graphs are execution objects, not semantic edges.\\
Value-aware block utility & No & Native-projection numerical patches do not evaluate that utility.\\
Budgeted submodular selection & No & A source-slot convention is not an evaluated submodular selector.\\
Exact/prototype/residual tiers & No & Exact overlays and sparse arithmetic patches are a different format.\\
Uncertainty-guided residual repair & No & History reconstruction is unconditional, not a repair policy.\\
\bottomrule
\end{tabularx}
\caption{Explicit scope of the current variant. Existing modules, CPU tests
or historical diagrams cannot turn an inactive mechanism into measured
evidence. These mechanisms remain research directions, not contributions
validated by the current performance table.}
\end{table*}

\section{Baseline provenance and historical evidence}
The current labels distinguish FullKV reference, Prompt Top-K adaptation,
and the PrefixKV, VL-Cache and LOOK-M reimplementation-plus-adaptation
paths. Their paper citations identify algorithmic origins, not ownership
of the present code or a claim that those papers used this codec
\citep{wang2025prefixkv,tu2025vlcache,wan2024lookm}.

The repository also preserves separate historical studies: native
PrefixKV Description/MM-Vet PPL and ROUGE-L routes; shared-byte FP16
comparisons; LLaVA-v1.6-Mistral COCO/CIDEr runs; and representative
MileBench Accuracy/ROUGE-L evaluations. They differ in checkpoints,
precision, outputs, budgets, variants or dataset routes. Their quality
numbers cannot be pooled with this BF16 LLaVA-1.5 storage study, and their
official-path checks do not automatically validate this new adaptation.
The main result table does not import those historical values.

In particular, a published result is neither a measured local row nor an
automatic comparator at the same actual bytes. Reproduction deltas require
their own identical native controls. The general lower-memory-at-equal-quality
question requires budget curves and an independently frozen eligibility
rule; comparing the current isolated operating points cannot answer it.

\section{Artifact verification and source map}
The paper exporter reads the immutable plan, source/ledger seals, every
worker manifest/result/raw record, phase decisions and the completed
statistical summary. It checks 124 producer workers, unique row keys,
complete sample grids, no FullKV fallback, identical paired tokens,
storage endpoints and source revisions. It independently recalculates
memory means/maxima and the timing mean/p50/p95. It binds the resulting
tables and plot data to SHA-256 hashes. Statistical intervals are exported
from the completed corpus/bootstrap analysis, not fabricated from rounded
table entries.

The exact producer is listed in Table~\ref{tab:environment}. Relevant
modules under \texttt{src/mosaickv/} include:
\begin{itemize}
\item \smallpath{evaluation/common_streaming_timing.py}: common explicit
      request lifecycle and measurement path;
\item \smallpath{evaluation/streaming_source_constructor.py}: layer-local
      source construction during the one prefill;
\item \smallpath{evaluation/shared_source_format.py}: persistent format
      and reconstruction contract;
\item the precision-tier, direct-text-construction, source-history and
      native-history-reader modules: payload coding, exact overlays,
      history sealing and bitwise checks;
\item the uniform-eager-core and transport modules: layout eligibility,
      request-private eager replay and retirement.
\end{itemize}
This map describes the measured immutable worktree. Unrelated newer or
dirty source edits are not implied to have produced the results.

Earlier fixture/schema failures remain separate from the accepted third
attempt. A scheduler completion code alone is insufficient: the exporter
also requires successful result decisions and complete raw denominators.
No unfavorable accepted row is dropped, and no GPU work is performed by
the document exporter or TeX build.

\section{Review-derived concern map and remaining work}
The supplied review-derived instructions requested complete-path timing,
physical ownership accounting, independent transport checks, faithful
baseline comparisons, boundary controls, mechanism ablations, fresh
confirmation and metric-aware uncertainty. The current paper addresses
some measurement concerns within a changed scope and explicitly withdraws
claims that would require the other experiments.

A separate private response gives a point-by-point account without
pretending that the unavailable original assessment is quoted verbatim.
No reviewer identity is required in the anonymous paper. Withdrawal of a
claim is not completion of the experiment that would have justified it.

Before broader conclusions, the remaining work is to match measured bytes,
certify whole-request cache peaks, validate native baseline fidelity for
the present adaptation, measure additional models/tasks and independently
confirm quality. For the historical architecture, each of its six active
components would additionally require correctness and causal-ablation
evidence. Multi-request serving, long contexts, GQA, lower-bit portability,
batch scaling and energy remain unmeasured opportunities, not established
benefits.

\clearpage
\bibliographystyle{mlsys2025}
\bibliography{references}
\end{document}
